\section{Limitations}
\label{sec:limitations}

Three limitations bound the claims in this paper.

\textbf{Evaluation scope.} Our study covers only machine learning papers, where computational reproducibility and well-defined contribution types make \ara{}'s four-layer structure a natural fit; whether the protocol generalizes to experimental sciences with physical execution requirements, or to theoretical disciplines where the Physical Layer is largely absent, remains empirically untested. Extending the Physical Layer to formal or proof-based results requiring machine-checkable specifications is a natural direction for future work. Our human-annotated benchmark was also constructed by annotators familiar with both the \ara{} format and the selected papers, so performance on unfamiliar or niche-domain artifacts may differ from the reported figures.

\textbf{Fidelity ceiling.} \ara{} fidelity is bounded by the source of supervision. The Compiler faithfully represents only what the PDF contains (\S\ref{sec:ingestor}): when a paper omits experimental details, environment specifications, or ablation results, no extraction method can recover them. The Live Research Manager closes this gap by recording trajectories as research unfolds, but assumes an AI-native workflow in which a coding agent is present throughout the project. For researchers outside such sessions, the Compiler still produces a valid \ara{} from the finished paper, but the resulting artifact inherits the PDF's omissions; hand-authoring structured fields remains possible but reintroduces the documentation burden the protocol aims to eliminate. Closing this adoption gap tracks the broader diffusion of agent workflows in research practice, not the protocol itself.

\textbf{Deployment prerequisites.} Two properties required for production use are not yet implemented. The adversarial robustness and privacy guarantees raised in \S\ref{sec:liveness} are aspirational: the current system lacks sandboxed execution, content-level anomaly detection, and granular access control for the Exploration Graph. Separately, any long-lived format faces \emph{schema evolution}: as research practice changes, the \ara{} schema will need to add node types, refine field semantics, and deprecate conventions without breaking prior artifacts. We version \texttt{PAPER.md} frontmatter with an explicit \texttt{ara\_schema} tag and require all validators to accept unknown fields (forward compatibility) and degrade gracefully on missing optional fields (backward compatibility), but have only exercised this discipline across minor revisions. A stable migration story for major revisions, including automatic rewriting of archival artifacts, long-term checker availability, and a deprecation policy, remains future work.
